\subsection{Calibration}\label{\refsec:calibration}
The calibration is a nested fixed point. In the inner part we solve the full politico-economic equilibrium path of \cref{\refsec:PEEpath}; in the outer part we search over the four parameters that cannot be set without it. The target system is repeated from \cref{\refmodelsec:calibration}:
\begin{subequations}\label{\refeq:calibration}
    \begin{align}
        \hat{\text{sr}}_{t_0} &= \frac{s_{t_0}}{\left(s_{t_0-1}/\nu_{t_0}\right)^{\alpha} h_{t_0}^{1-\alpha}} \label{\refeq:calibration:sr}\\
        \hat{\tau}_{t_0} &= \tau_{t_0} \label{\refeq:calibration:tau}\\
        \hat{\eta}_{0} &= \frac{z_0^{\eta}}{z_0^x}\frac{(1-\alpha)(1-\tau_{t_0})}{\Theta_{h,t_0}^{\alpha}\left(\frac{1-\alpha}{\Gamma_{h,t_0}^{\alpha}}\right)^{\frac{1}{1+\alpha\xi}}} \label{\refeq:calibration:eta0}\\
        \hat{X}_{0} &= \hat{\eta}_{0}\frac{\left(\frac{1-\alpha}{\Gamma_{h,t_0}^{\alpha}}\right)^{\frac{1}{1+\alpha\xi}}}{\left(\Theta_{h,t_0}z_0^x\right)^{1/\xi}} \label{\refeq:calibration:X0}
    \end{align}
\end{subequations}
where the ``hats'' reflect exogenous targets and $t_0$ indicates a specific baseline year. In our main specification $\beta_{t,j} = \beta$ is constant over time and across types, and most exogenous parameters are constant over time except for the ageing population reflected by $\nu_t$. Loosely, $\beta$ carries the savings rate, $\omega$ the tax rate, and $\eta_0, X_0$ the informal income and labour targets --- but as we note below, that reading of the Jacobian describes the dominant terms only and should not be relied on.

Everything else is set before the loop and does not depend on it: $\bm{\eta}_i,\bm{X}_i$ from the eigenvector step \eqref{\refmodeleq:calibration:etai}--\eqref{\refmodeleq:calibration:Xi}, then $\theta$ from the replacement-rate ratio and $\epsilon$ from the coverage rate. Note that $\epsilon$ does depend on $\beta$, and $\kappa$ in turn on $\epsilon$, so these two auxiliary parameters must be refreshed inside the loop even though the expressions defining them are not part of \eqref{\refeq:calibration}.

\smalltitle{Two parameters that are not calibrated.}
The savings rate \eqref{\refeq:calibration:sr} deliberately measures the \emph{formal} sector only: $s_{t_0}$ excludes the informal households' savings $\iota_{t_0}s_{t_0}$ per formal young household, and the denominator excludes informal output $w_{t_0}^0\eta_{t_0,0}h_{t_0,0}$ correspondingly. Numerator and denominator therefore cover the same population, and $\iota$ is left free to serve as a diagnostic rather than as something the calibration has been fitted to. This is a choice, not an oversight: a national-accounts reading of the target would put $s_{t_0}\left(1+\gamma_{t_0,0}\iota_{t_0}\right)$ over total output instead, and would make $\beta$ answerable to the informal block as well.

The informal return scaler $\chi_t^R$ is likewise not calibrated. It is fixed outside the loop, alongside $\xi$ and $\rho$, because nothing in the data we use identifies it. It is worth being explicit that $\chi^R = 1$ is not a neutral default: it asserts that the informal savings vehicle earns exactly the formal return, which is the knife-edge at which the vehicle differs from formal saving only through the pension system and not through its return. The natural reading is $\chi^R\leq 1$, and results should be reported across a range of $\chi^R$ rather than at a single value. Note that $\chi^R$ reaches the equilibrium through $B_{t+1}^0$ in \eqref{\refeq:auxiliary:B0} and through $A_t$ in \eqref{\refeq:auxiliary:Ainf}, hence through both the level and the political weight of the informal savings ratio; it is not a re-scaling that washes out.

\smalltitle{Why $\eta_0$ and $X_0$ are the awkward pair.}
The informal productivity and disutility parameters reach the equilibrium through two distinct channels. They set the informal households' consumption levels, and hence the weight attached to the informal terms of the first order condition; and, by \eqref{\refeq:auxiliary:s0_s}, the ratio $\iota_t$ is proportional to $\eta_{t,0}^{1+\xi}/X_{t,0}^{\xi}$, so they also set the scale of the political state itself. A trial $(\eta_0,X_0)$ therefore shifts the whole region of the state space in which the equilibrium lives, not merely a coefficient inside it --- and it does so before a single grid has been laid down.

Three consequences follow, and all three are practical rather than conceptual.
\begin{itemize}
    \item \emph{The state grids are parameter-dependent and must be rebuilt at every residual evaluation.} The bounds $l_{\iota},u_{\iota}$ of \cref{\refsec:PEELOG} are constructed from the steady-state map $\iota^*(\tau)$, which moves with $(\beta,\eta_0,X_0)$. A grid built once at the starting parameters and reused across the search would silently place the equilibrium outside its own grid, and the failure would show up as an infeasibility or a corner solution rather than as an error. The same applies to $\mathcal{S}$ in the CRRA case, whose bounds come from $s^*(\tau)$ and therefore move with $\beta$.
    \item \emph{The Jacobian is not close to diagonal.} $\eta_0$ and $X_0$ move $\tau_{t_0}$ through the state channel, and $\Theta_{h,t_0}$ appears on the right-hand side of both \eqref{\refeq:calibration:eta0} and \eqref{\refeq:calibration:X0}, so the two self-consistency conditions are coupled to the tax target as well as to each other. The parameter-to-target assignment above should be read as a description of the dominant terms, not as a basis for solving the four equations sequentially.
    \item \emph{The outer search needs a starting point, and there is a cheap one.} Freezing $\iota_{t-1}$ at the steady-state value $\iota^*(\tau)$ instead of tracking it as a state collapses the political problem to a scalar problem per period, of the kind that can be solved for the whole path at once. The parameters that calibration returns are not the ones we want --- the frozen state is exactly the feedback we are interested in --- but they are obtained at a small fraction of the cost and are a far better starting point for the full search than an arbitrary guess.
\end{itemize}

\smalltitle{The outer residual is only piecewise smooth.}
This is the main numerical hazard of the whole procedure, and it comes from the inner solve passing through a grid search: the composed map from parameters to targets inherits two kinds of non-smoothness, and neither is visible in \eqref{\refeq:calibration} itself.

Within a branch it is better behaved than one might fear: the selection rule \eqref{\refeq:candidates} locates $\tau_t$ by linear interpolation \emph{inside} a cell rather than snapping to a node, so a small parameter perturbation moves $\tau_t$ continuously. What it does not do is remove kinks. The policy interpolants $\tau^t(\cdot)$ --- and $\tau^t(\cdot,\cdot)$ in the CRRA case --- are piecewise linear in the state, so the map is continuous but only piecewise continuously differentiable, with kinks spaced at the resolution of the state grid. And it can jump outright, when the $\arg\max$ in \eqref{\refeq:candidates} switches between well-separated candidates, or when the feasibility mask gains or loses a node.

It is natural to conclude from this that the finite-difference step used to build the outer Jacobian must be large enough for the induced movement of the state to clear a grid cell, so that the difference quotient measures the economics rather than the interpolation, and that a step chosen by the usual $\sqrt{\text{machine precision}}$ heuristic will be far below that and return noise. \emph{That conclusion is wrong, and in the log case the prescription it implies is actively harmful.} The reason is in the previous paragraph: because $\tau_t$ is located by interpolation \emph{inside} a cell and the policy interpolants are piecewise linear in the state, a small step stays on one linear piece and returns that piece's slope --- which is precisely what a Newton step wants --- whereas a large step averages across kinks. Measured at the converged log calibration, the difference quotient is flat across six orders of magnitude of step size and it is the \emph{large} steps that move:
\begin{center}
\begin{tabular}{lccccc}
\toprule
step $h$ & $10^{-9}$ & $1.49\cdot 10^{-8}$ & $10^{-5}$ & $10^{-3}$ & $10^{-2}$\\
\midrule
$\partial(\text{res}_{\tau})/\partial(\ln\beta)$ & $-0.1352$ & $-0.1352$ & $-0.1352$ & $-0.1096$ & $-0.1287$\\
$\partial(\text{res}_{\tau})/\partial(\ln\omega)$ & $+0.3698$ & $+0.3698$ & $+0.3697$ & $+0.3826$ & $+0.3837$\\
\bottomrule
\end{tabular}
\end{center}
where $1.49\cdot 10^{-8}=\sqrt{\text{machine precision}}$ is the default. Agreement between $h=10^{-9}$ and the default step is $9\cdot 10^{-5}$ relative or better on all four columns, and the log calibration converges from arbitrary starting parameters in $26$ evaluations at that step. In the log case the hazard is therefore real but \emph{intermittent} --- a step that happens to straddle a kink, or an $\arg\max$ that switches --- rather than a systematic noise floor to be engineered around.

\smalltitle{A measurement that appeared to say the opposite, and did not.} The paragraph above was measured at a converged \emph{log} calibration, and it is worth recording what happened when the same measurement was repeated in the CRRA case, because the episode is more instructive than the number it produced. At $\rho=1.1$ on the $45\times 45$ grid, starting from the log parameters, three of the four columns behaved and one did not:
\begin{center}
\begin{tabular}{lcccccc}
\toprule
step $h$ & $1.49\cdot 10^{-8}$ & $10^{-6}$ & $10^{-5}$ & $10^{-4}$ & $10^{-3}$ & $10^{-2}$\\
\midrule
$\beta$   & $0.132$ & $0.132$ & $0.131$ & $0.138$ & $0.137$ & $0.136$\\
$\omega$  & $0.386$ & $0.451$ & $0.350$ & $0.398$ & $0.387$ & $0.399$\\
$\eta_0$  & $\mathbf{5.126}$ & $0.551$ & $0.905$ & $\mathbf{0.988}$ & $\mathbf{0.999}$ & $0.916$\\
$X_0$     & $0.996$ & $0.996$ & $0.996$ & $0.997$ & $0.992$ & $0.990$\\
\bottomrule
\end{tabular}
\end{center}
$\beta$ and $X_0$ are stable at every step and $\omega$ wobbles by $\pm 13\%$, but the $\eta_0$ column at the default step is \emph{five times} its resolved value, and a single corrupted column is enough to spoil the Newton direction. The reading at the time was that the log conclusion simply does not generalise, and the CRRA outer solve was run at $h=10^{-4}$ on that basis.

That reading was wrong, and the cause was not in the step at all. The policy smoother then in use chose its knot \emph{count} from the data, which put small jumps in the outer residual (\cref{\refsec:PEELOG}); $h=1.49\cdot 10^{-8}$ was simply a step that straddled one. Once the smoother's knots are pinned the column is flat --- within $0.01\%$ of its value at $h=10^{-4}$ over the whole range $[1.5\cdot 10^{-8}, 10^{-4}]$, and within $0.6\%$ even at $10^{-2}$ --- at both of the converged points it was re-measured at. Calibrating from a common start at each candidate step then converges in the same number of evaluations to the same parameters to ten digits, with the \emph{default} step reaching the tightest final residual of the three. Both preference cases therefore run at $\sqrt{\text{machine precision}}$, and the first paragraph's conclusion holds generally rather than case by case.

The point worth carrying away is not about step sizes. A diagnostic taken at one point of one case is a statement about that residual, and if the residual has a defect elsewhere in it, the diagnostic will faithfully report the defect as a property of the thing being measured. Both readings above were correct measurements; only one of them was a correct inference.

\smalltitle{Where the hazard does bite: the inner grid, not the step.}
In the CRRA case the composed map is rough enough to matter, and the object at fault is the resolution of the inner state grids. At the resolution the politico-economic solve itself uses ($\vert\mathcal{S}\vert=\vert\mathcal{S}_0\vert=30$), the outer search converges --- and reports success --- to a \emph{displaced} point. The diagnostic is to hold the converged parameters fixed and refine only the inner grid:
\begin{center}
\begin{tabular}{lccccc}
\toprule
$(\vert\mathcal{S}_0\vert,\vert\mathcal{S}\vert,\text{refinement of }\mathcal{S}')$ & $(30,30,4)$ & $(30,30,8)$ & $(45,45,4)$ & $(60,60,4)$ & $(60,60,8)$\\
\midrule
$\max\vert\text{residual}\vert$ & $6.6\cdot 10^{-4}$ & $7.6\cdot 10^{-4}$ & $3.13\cdot 10^{-3}$ & $3.14\cdot 10^{-3}$ & $3.14\cdot 10^{-3}$\\
\bottomrule
\end{tabular}
\end{center}
The sequence \emph{plateaus} rather than decaying, which says the refined problem is well resolved and its root simply lies elsewhere: on that evidence the coarse answer was wrong, not imprecise. Calibrating at $45\times 45$ instead gave the pattern one wants --- $10^{-12}$ on its own grid, and $1.0\cdot 10^{-4}$ and $4.4\cdot 10^{-5}$ when re-evaluated at $60\times 60$ and $75\times 75$ --- and $45\times 45$ was adopted on that basis. The log case needs no such treatment: at its own default its residual under refinement already shrinks.

This measurement, like the step-size one above it, was taken before the smoother's knots were pinned, and it too weakens considerably when repeated afterwards. The displacement largely disappears: re-calibrating at each grid size and comparing the converged \emph{parameters} rather than the residual, all of $\vert\mathcal{S}_0\vert=\vert\mathcal{S}\vert\in\{30,45,60\}$ now agree to between $1.5\cdot 10^{-4}$ and $3.8\cdot 10^{-4}$ in relative terms, against the roughly $1\%$ that the plateau above was reporting. What remains is a genuine but modest resolution effect, and only at the harder end of the parameter range: the refinement rungs at $30\times 30$ sit $2$--$3$ times above those at $45\times 45$ when $\rho=0.7$, and are indistinguishable from them when $\rho=0.9$. Going beyond $45\times 45$ buys nothing measurable --- $60\times 60$ produces the same rung levels and parameters differing by $10^{-4}$, the same order as the spread across all three --- so about $10^{-4}$ in the parameters is the floor of what the outer answer is determined to, and is the scale against which the outer tolerance should be read. We keep $45\times 45$, at roughly four times the cost of $30\times 30$, on the narrow ground that the sweep extends to $\rho = 0.5$, further into the region where the gap opened at all.

The general lesson is worth stating separately, because both of the usual reassurances were present at the wrong answer: the root finder reported success, and the residual was small. What distinguished the two cases was refining the grid the residual is computed \emph{through} and asking whether the number decays or plateaus --- one extra evaluation, and nothing about the outer solve's own tolerance could have substituted for it.

\smalltitle{Deeper than the resolution: the \emph{kind} of continuation interpolant.}
Refining the state grids treats a symptom. The cause is named in the second paragraph of this section and then not acted upon: the policy interpolants are piecewise linear in the state, so the composed map is continuous but only piecewise continuously differentiable. Grid refinement moves the kinks closer together without removing them, and a Newton step still cannot descend a surface that is kinked at every cell boundary. Replacing the piecewise-linear continuation surfaces with $C^1$ ones removes the kinks outright. Measured at $\rho=1.1$, $45\times 45$, at the $h=10^{-4}$ step then in use:
\begin{center}
\begin{tabular}{lcc}
\toprule
 & piecewise linear & cubic\\
\midrule
outer solve & stalls at $\approx 3\cdot 10^{-5}$, $60+$ evaluations & $1.33\cdot 10^{-11}$ in $12$ evaluations\\
residual at the converged point, $45/60/75$ & $3.3\cdot 10^{-5}\to 1.5\cdot 10^{-4}\to 2.9\cdot 10^{-4}$ & $1.3\cdot 10^{-11}\to 1.7\cdot 10^{-4}\to 5.1\cdot 10^{-5}$\\
\bottomrule
\end{tabular}
\end{center}
The refinement trend is the diagnostic of the previous paragraph applied to itself, and under piecewise-linear interpolation it has the \emph{wrong sign}: the residual grows, so that answer is not grid-converged at any of the three resolutions. Under cubic interpolation it decays, which is the pattern the previous paragraph calls healthy. The calibrated parameters differ between the two by only $0.016\%$ in $\beta$ and $0.11\%$ in $\omega$, so the piecewise-linear answer was not badly \emph{displaced} --- the solver simply could not drive its residual down. The CRRA calibration is therefore run with cubic continuation interpolants.

A monotone $C^1$ interpolant (PCHIP) would be preferable to a plain cubic, which can overshoot where a policy is flat at a bound --- on a test profile flat at both ends of $[0,3]$, cubic returns values in $[-0.088,\,3.105]$ where PCHIP stays inside. The tax rate is clipped back into $[l,u]$ in any case (see the smoothing discussion), but the remaining surfaces are not. This is not currently done for a purely computational reason: the available implementation rebuilds its splines at every evaluation, costing some three orders of magnitude more than the cubic, which would dominate the entire solve.

Two further points on robustness. An inner solve that reports infeasibility --- a path leaving the reachable set \eqref{\refeq:reachable}, or the initial fixed point \eqref{\refeq:initialFixedPoint} failing to bracket --- should propagate that failure rather than return a large sentinel residual, since a sentinel invites the root finder to wander into a region where every evaluation fails identically. And the parameter set should never be left holding a trial point: a failed outer iteration must restore the parameters that were in place when the calibration started, so that a later diagnostic is not run against a state the calibration merely passed through.

\begin{alg}[Nested fixed point calibration]\label{\refalg:calibration}
Fix $\xi,\rho,\chi^R$ and the targets $\hat{\text{sr}}_{t_0},\hat{\tau}_{t_0},z_0^{\eta},z_0^x$. Set $\bm{\eta}_i,\bm{X}_i$ and $\theta$ once, outside the loop.

\smallskip\noindent\emph{Outer step.} Search over $\left(\beta,\omega,\eta_0,X_0\right)$, each mapped to the real line by a logarithm so that positivity is enforced without a bound being reached. Note $\beta$ is \emph{not} capped at unity: with $\beta_{t,j} = \beta p_{t,j}$ the actual discount factor is already below one, and imposing $\beta<1$ makes the search stall on the constraint rather than fail informatively.

\smallskip\noindent\emph{Inner step.} Given a trial $\left(\beta,\omega,\eta_0,X_0\right)$:
\begin{enumerate}
    \item Write the parameters into the parameter set and refresh the auxiliary definitions that depend on them, in order: $\epsilon$ from $\beta$, then $\kappa$ from $\epsilon$, then $\Gamma_h$ from $\left(\bm{\eta},\bm{X}\right)$.
    \item Rebuild the state grids at the trial parameters --- $\mathcal{S}_0$ from $\iota^*(\tau)$, and $\mathcal{S}$ from $s^*(\tau)$ in the CRRA case --- as discussed above. Never carry them over from the previous evaluation. In the CRRA case rebuild them at $\vert\mathcal{S}\vert=\vert\mathcal{S}_0\vert=45$ rather than at the resolution the politico-economic solve alone would use, and build the continuation interpolants as cubic rather than piecewise-linear surfaces; both for the reasons measured above.
    \item Solve the politico-economic equilibrium: \cref{\refalg:LOG:gridsearch} or \cref{\refalg:CRRA:grid} for the policy functions, then \cref{\refsec:PEEpath} for the initial state, the forward simulation and the exact re-solve of the economic equilibrium at the simulated $\bm{\tau}$.
    \item Read off the four quantities in \eqref{\refeq:calibration} at $t_0$, recovering $\Theta_{h,t_0}$ from the solved $\left(h_{t_0},s_{t_0-1}\right)$ by inverting \eqref{\refeq:auxiliary:h}, which covers the log and CRRA cases and the general and terminal formulas alike with one expression.
\end{enumerate}

\smallskip\noindent\emph{Residual.} Form the four gaps with \eqref{\refeq:calibration:sr} and \eqref{\refeq:calibration:tau} as level differences --- both are rates of the same order of magnitude --- and \eqref{\refeq:calibration:eta0} and \eqref{\refeq:calibration:X0} as relative differences. At the calibrated point all four quantities turn out to be of the same order ($\text{sr}=0.184$, $\tau=0.125$, $\eta_0=0.326$, $X_0=0.408$), so the choice of scaling changes the last two rows only by a factor of two or three and is close to immaterial; it is kept because a relative gap remains the natural form for a self-consistency condition, not because the search needs it.

\smallskip\noindent\emph{Outer step size.} Take the outer Jacobian at the default $\sqrt{\text{machine precision}}$ step in the log case and at $h=10^{-4}$ in the CRRA case, for the reason measured above. This is the one setting that cannot be left to the root finder's defaults.
\end{alg}

\smalltitle{Order of work, and a check.}
The log case is the cheaper of the two --- one state, and by \eqref{\refeq:zdecomposition} a state that costs almost nothing to carry --- so it is the natural place to calibrate first, and its converged parameters are the natural starting point for the CRRA case. Continuing from $\rho = 1$ outwards in small steps is worthwhile for a second reason: the two solvers share no code path, so agreement between the log calibration and a CRRA calibration at $\rho$ close to $1$ is a genuine test of both, and it is a test that only works if the CRRA run is started close enough to converge at all. Once calibrated, $\iota_{t_0}$ is worth inspecting even though nothing was fitted to it: it is the object the informal savings decision exists to produce, and the whole grid construction rests on the restriction $l_{\iota}>0$, which a calibrated value outside the unit interval would contradict.

\smalltitle{Calibrating across a grid of $\rho$.}
Results are wanted not at one $\rho$ but across a range of it, and a calibration costs a few dozen politico-economic solves, so the sequence is worth organising rather than looping over. Three points, in order of how much they matter.

First, the sequence is \emph{anchored}, not swept: $\rho=1$ is solved first because it is the only value at which the cheaper log solver applies and the only one needing no starting point at all, and the grid is then walked outwards in both directions from it, each direction carrying its own history. Second, each calibration is started from an extrapolation through the two already solved nearest it, taken in the \emph{unbounded} coordinate of the outer search rather than in the parameters themselves --- extrapolating a positivity-constrained parameter in levels can step across its own bound, and the resulting starting point is not merely poor but invalid. Third, a value that refuses to solve is retried from the un-extrapolated previous point and then, failing that, approached through an intermediate value inserted between it and the last success; the intermediate is kept, so the retry extrapolates from closer in.

Two observations from running this. The march's own diagnostic is $\eta_0$ and $X_0$: across $\rho\in[0.9,1.1]$ they move by $0.00\%$ and $0.02\%$ respectively, against $29\%$ for $\beta$ and $22\%$ for $\omega$ --- so the two self-consistency conditions are very nearly solved by direct substitution, and it is only $\left(\beta,\omega\right)$ that the data targets are genuinely identifying. And the extrapolation earns less than one would expect once the interpolants are $C^1$: at a step of $\Delta\rho=0.1$ a warm-started calibration and a cold one both converge in twelve evaluations, to identical parameters. Its value is in robustness at larger steps and in surviving a failure, not in speed.
